% !TEX TS-program = xelatex
% !TeX spellcheck = ru_RU
% Файл можно включить в диплом командой % !TEX TS-program = xelatex
% !TeX spellcheck = ru_RU
% Файл можно включить в диплом командой % !TEX TS-program = xelatex
% !TeX spellcheck = ru_RU
% Файл можно включить в диплом командой % !TEX TS-program = xelatex
% !TeX spellcheck = ru_RU
% Файл можно включить в диплом командой \input{task}.
% При самостоятельной компиляции используется оформление предпросмотра.
\ifdefined\maketitle
    \let\taskdocumentend\relax
\else
    \documentclass[12pt,a4paper]{article}
    \usepackage{fontspec}
    \setmainfont{cmunrm.otf}[BoldFont=cmunbx.otf,ItalicFont=cmunti.otf,BoldItalicFont=cmunbi.otf]
    \usepackage{polyglossia}
    \setdefaultlanguage{russian}
    \usepackage[top=20mm,bottom=20mm,left=30mm,right=15mm]{geometry}
    \usepackage{setspace}
    \onehalfspacing
    \setlength{\emergencystretch}{2em}
    \def\taskdocumentend{\end{document}}
    \begin{document}
\fi

\section{Постановка задачи}
\label{sec:task}

Целью работы является разработка и экспериментальная оценка
способов распределения ограниченной памяти между дообучением модели
и калибровкой порога сетевого детектора при возвратном дрейфе
нормального трафика.

Для её выполнения были поставлены следующие задачи:
\begin{enumerate}
    \item Подготовить сравнительный обзор методов адаптации сетевых
          детекторов, сохранения исторических данных и калибровки
          порогов, выделить подходы к распределению ограниченной памяти
    \item Сформулировать задачу распределения общего бюджета памяти
          между дообучением и калибровкой порога и определить критерии
          оценки результата
    \item Разработать способы определения долей наблюдений,
          сохраняемых для дообучения модели и калибровки её порога
    \item Реализовать предложенные способы и базовые варианты
          с фиксированными долями памяти для выбранного сетевого детектора
    \item Подготовить воспроизводимые сценарии изменения и возвращения
          режимов нормального трафика на сетевых данных, содержащих
          нормальные наблюдения и атаки
    \item Экспериментально сравнить предложенные способы с базовыми
          вариантами при одинаковом бюджете памяти, оценить влияние
          распределения памяти на долю ложных тревог при возвращении
          нормальных режимов и долю обнаруженных атак, определить
          условия полезности способов и их ограничения
\end{enumerate}

\taskdocumentend
.
% При самостоятельной компиляции используется оформление предпросмотра.
\ifdefined\maketitle
    \let\taskdocumentend\relax
\else
    \documentclass[12pt,a4paper]{article}
    \usepackage{fontspec}
    \setmainfont{cmunrm.otf}[BoldFont=cmunbx.otf,ItalicFont=cmunti.otf,BoldItalicFont=cmunbi.otf]
    \usepackage{polyglossia}
    \setdefaultlanguage{russian}
    \usepackage[top=20mm,bottom=20mm,left=30mm,right=15mm]{geometry}
    \usepackage{setspace}
    \onehalfspacing
    \setlength{\emergencystretch}{2em}
    \def\taskdocumentend{\end{document}}
    \begin{document}
\fi

\section{Постановка задачи}
\label{sec:task}

Целью работы является разработка и экспериментальная оценка
способов распределения ограниченной памяти между дообучением модели
и калибровкой порога сетевого детектора при возвратном дрейфе
нормального трафика.

Для её выполнения были поставлены следующие задачи:
\begin{enumerate}
    \item Подготовить сравнительный обзор методов адаптации сетевых
          детекторов, сохранения исторических данных и калибровки
          порогов, выделить подходы к распределению ограниченной памяти
    \item Сформулировать задачу распределения общего бюджета памяти
          между дообучением и калибровкой порога и определить критерии
          оценки результата
    \item Разработать способы определения долей наблюдений,
          сохраняемых для дообучения модели и калибровки её порога
    \item Реализовать предложенные способы и базовые варианты
          с фиксированными долями памяти для выбранного сетевого детектора
    \item Подготовить воспроизводимые сценарии изменения и возвращения
          режимов нормального трафика на сетевых данных, содержащих
          нормальные наблюдения и атаки
    \item Экспериментально сравнить предложенные способы с базовыми
          вариантами при одинаковом бюджете памяти, оценить влияние
          распределения памяти на долю ложных тревог при возвращении
          нормальных режимов и долю обнаруженных атак, определить
          условия полезности способов и их ограничения
\end{enumerate}

\taskdocumentend
.
% При самостоятельной компиляции используется оформление предпросмотра.
\ifdefined\maketitle
    \let\taskdocumentend\relax
\else
    \documentclass[12pt,a4paper]{article}
    \usepackage{fontspec}
    \setmainfont{cmunrm.otf}[BoldFont=cmunbx.otf,ItalicFont=cmunti.otf,BoldItalicFont=cmunbi.otf]
    \usepackage{polyglossia}
    \setdefaultlanguage{russian}
    \usepackage[top=20mm,bottom=20mm,left=30mm,right=15mm]{geometry}
    \usepackage{setspace}
    \onehalfspacing
    \setlength{\emergencystretch}{2em}
    \def\taskdocumentend{\end{document}}
    \begin{document}
\fi

\section{Постановка задачи}
\label{sec:task}

Целью работы является разработка и экспериментальная оценка
способов распределения ограниченной памяти между дообучением модели
и калибровкой порога сетевого детектора при возвратном дрейфе
нормального трафика.

Для её выполнения были поставлены следующие задачи:
\begin{enumerate}
    \item Подготовить сравнительный обзор методов адаптации сетевых
          детекторов, сохранения исторических данных и калибровки
          порогов, выделить подходы к распределению ограниченной памяти
    \item Сформулировать задачу распределения общего бюджета памяти
          между дообучением и калибровкой порога и определить критерии
          оценки результата
    \item Разработать способы определения долей наблюдений,
          сохраняемых для дообучения модели и калибровки её порога
    \item Реализовать предложенные способы и базовые варианты
          с фиксированными долями памяти для выбранного сетевого детектора
    \item Подготовить воспроизводимые сценарии изменения и возвращения
          режимов нормального трафика на сетевых данных, содержащих
          нормальные наблюдения и атаки
    \item Экспериментально сравнить предложенные способы с базовыми
          вариантами при одинаковом бюджете памяти, оценить влияние
          распределения памяти на долю ложных тревог при возвращении
          нормальных режимов и долю обнаруженных атак, определить
          условия полезности способов и их ограничения
\end{enumerate}

\taskdocumentend
.
% При самостоятельной компиляции используется оформление предпросмотра.
\ifdefined\maketitle
    \let\taskdocumentend\relax
\else
    \documentclass[12pt,a4paper]{article}
    \usepackage{fontspec}
    \setmainfont{cmunrm.otf}[BoldFont=cmunbx.otf,ItalicFont=cmunti.otf,BoldItalicFont=cmunbi.otf]
    \usepackage{polyglossia}
    \setdefaultlanguage{russian}
    \usepackage[top=20mm,bottom=20mm,left=30mm,right=15mm]{geometry}
    \usepackage{setspace}
    \onehalfspacing
    \setlength{\emergencystretch}{2em}
    \def\taskdocumentend{\end{document}}
    \begin{document}
\fi

\section{Постановка задачи}
\label{sec:task}

Целью работы является разработка и экспериментальная оценка
способов распределения ограниченной памяти между дообучением модели
и калибровкой порога сетевого детектора при возвратном дрейфе
нормального трафика.

Для её выполнения были поставлены следующие задачи:
\begin{enumerate}
    \item Подготовить сравнительный обзор методов адаптации сетевых
          детекторов, сохранения исторических данных и калибровки
          порогов, выделить подходы к распределению ограниченной памяти
    \item Сформулировать задачу распределения общего бюджета памяти
          между дообучением и калибровкой порога и определить критерии
          оценки результата
    \item Разработать способы определения долей наблюдений,
          сохраняемых для дообучения модели и калибровки её порога
    \item Реализовать предложенные способы и базовые варианты
          с фиксированными долями памяти для выбранного сетевого детектора
    \item Подготовить воспроизводимые сценарии изменения и возвращения
          режимов нормального трафика на сетевых данных, содержащих
          нормальные наблюдения и атаки
    \item Экспериментально сравнить предложенные способы с базовыми
          вариантами при одинаковом бюджете памяти, оценить влияние
          распределения памяти на долю ложных тревог при возвращении
          нормальных режимов и долю обнаруженных атак, определить
          условия полезности способов и их ограничения
\end{enumerate}

\taskdocumentend
